\section{BT1092--BT1093 Explicit Quotient and Odd--Even Matter Cube}
\label{sec:bt1092-bt1093-explicit-quotient-cube}

The reservoir intertwiner now has an explicit quotient matrix.  Let
\[
F_{13}=\operatorname{span}(e_0,e_1,\ldots,e_{12})
\]
be the fixed/gauge-perp block and let
\[
\nu=e_0+e_1+\cdots+e_{12}
\]
be its scalar trace line.  Define
\[
\pi_{12}:F_{13}\to\mathbb C^{12},
\qquad
\pi_{12}(x_0,\ldots,x_{12})=(x_0-x_{12},\ldots,x_{11}-x_{12}).
\]
Then
\[
\operatorname{rank}\pi_{12}=12,
\qquad
\ker\pi_{12}=\mathbb C\nu.
\]
Thus the BT1090 reservoir map is the concrete trace-free quotient
\[
K=\frac13\sum_{g=0}^{2}\pi_{12}^{(g)}\operatorname{pr}_{F_{13}(g)}.
\]

The $13$ itself is fixed by the antipodal decomposition of the matter cube.  With
\[
C^{27}_R(g)=\mathbb C[\mathbb F_3^3],
\]
the involution $x\mapsto -x$ fixes the origin and partitions the remaining $26$ vectors into the thirteen projective directions of $PG(2,3)$.  Hence
\[
\mathbb C[\mathbb F_3^3]
= B_{13}\oplus R_{14},
\]
where
\[
B_{13}=\operatorname{span}\{e_x-e_{-x}\},
\qquad
R_{14}=\mathbb C e_0\oplus\operatorname{span}\{e_x+e_{-x}\}.
\]
This gives
\[
27=13_{\mathrm{odd}}+14_{\mathrm{even}}.
\]
The right-chirality header is therefore the antipodal-odd matter-cube subspace, and the residual buffer is the antipodal-even subspace.
